\chapter{Post-Trade Systems}\label{ch:pl:post-trade-systems}

On 28 May 2024 most US securities trades began to settle one business day after the trade date instead of two. The rule that made the change
also required brokers to have agreements or policies aimed at completing allocations, confirmations and \omterm{def:pl:post-trade-systems:confirmation}{affirmations} by the end of the trade
date, and the industry's
plan moved the \omterm{def:pl:post-trade-systems:confirmation}{affirmation} cut-off of the central matching service from 11:30~a.m.\ on the day after the trade to 9:00~p.m.\ on the trade date
itself. For the operations team that fixes a trade whose confirmation does not agree with the firm's booking, the time available went from a
working day to an evening. This chapter builds the systems that decide how many trades need fixing --- matching, \omterm{def:pl:post-trade-systems:ssi}{settlement instructions},
\omterm{def:pl:post-trade-systems:recon}{reconciliation} --- and measures what the shorter cycle does to the ones that remain.

\section{From booked trade to settled trade}

Chapter~\ref{ch:pl:trade-capture-and-booking} ended with a booked trade in the canonical model. Between booking and settlement (Book~1,
chapter~5) the trade must be agreed with the counterparty, instructed to the agents who move the securities and cash, settled, and then checked
against what the custodian actually holds. Every step compares two records of the same thing, and every comparison that fails is a piece of
manual work with a deadline.

\begin{definition}[Trade confirmation, affirmation]\label{def:pl:post-trade-systems:confirmation}
A \emph{trade confirmation}\index{trade confirmation} is the executing party's statement of a trade's terms --- security, side, quantity, price,
trade and settlement dates, accounts --- sent to its counterparty. \emph{Affirmation}\index{affirmation} is the counterparty's (for an
institutional trade, the investment manager's or its custodian's) agreement that the confirmation is right, after which the trade can be
settled without further intervention.
\end{definition}

\begin{definition}[Settlement instruction, standard settlement instruction]\label{def:pl:post-trade-systems:ssi}
A \emph{settlement instruction}\index{settlement instruction} tells a settlement agent to deliver or receive a security against payment, with
the quantity, amount, settlement date and both parties' accounts. A \emph{standard settlement instruction}\index{standard settlement
instruction} (SSI) is the account details a counterparty uses for a given market, kept once as reference data so that each trade's
instruction is generated rather than typed.
\end{definition}

\begin{omfigure}
\begin{tikzpicture}[font=\small,node distance=4mm,
  box/.style={draw=omDef,fill=omDef!8,rounded corners=2pt,align=center,minimum height=10mm,text width=20mm},
  chk/.style={draw=omThm,fill=omThm!8,rounded corners=2pt,align=center,minimum height=10mm,text width=20mm},
  arr/.style={-{Stealth[length=2mm]},thick}]
\node[box] (b) {booked trade\\(ch.~21)};
\node[chk,right=of b] (m) {match with the\\confirmation};
\node[box,right=of m] (a) {affirm\\(cut-off)};
\node[box,right=of a] (i) {instruct\\(SSI table)};
\node[box,right=of i] (s) {settle\\(DVP at the CSD)};
\node[chk,below=9mm of s,text width=26mm] (r) {reconcile with\\the custodian};
\node[chk,below=9mm of m,text width=34mm] (q) {break queue: classify,\\assign, fix, age};
\draw[arr] (b) -- (m); \draw[arr] (m) -- (a); \draw[arr] (a) -- (i); \draw[arr] (i) -- (s); \draw[arr] (s) -- (r);
\draw[arr,omThm] (m) -- node[left,font=\scriptsize]{mismatch} (q);
\draw[arr,omThm] (r) -- node[above,font=\scriptsize]{difference} (q);
\draw[arr,omThm,dashed] (q.east) to[bend right=12] node[below right,font=\scriptsize,pos=0.35]{fixed} (a.south);
\end{tikzpicture}

\omcaption{The post-trade pipeline of this chapter: each booked trade is matched with the counterparty's confirmation, affirmed before the
cut-off, instructed from the SSI table and settled delivery versus payment; positions are then reconciled with the custodian. Every mismatch or
difference (red) becomes a break in one queue, where it is classified, owned, fixed and aged.}\label{fig:pl:post-trade-systems:pipeline}
\end{omfigure}

The chapter's week is a model, and every number in it comes from the model: 500 institutional equity trades a day, executed by eight brokers for
five funds, each trade booked in the firm's canonical model. The brokers' confirmations carry discrepancies planted at stated rates: a wrong
price for 1\% of trades, a wrong settlement date for 1\%, a wrong account for 1.5\%, a wrong quantity for 0.5\%, and no confirmation at all
for 1\%. Three of the eight brokers also round the average price to four decimals, which is not an error at all.

\section{Confirmation and affirmation}

\begin{definition}[Confirmation matching, matching tolerance]\label{def:pl:post-trade-systems:matching}
\emph{Confirmation matching}\index{confirmation matching} pairs each booked trade with a confirmation that agrees on key fields (security,
side, counterparty, trade date, account holder) and then compares the other fields; a \emph{matching tolerance}\index{matching tolerance} is
the largest difference in a compared field that still counts as agreement. A pair within every tolerance is matched; a pair outside some is a
partial match, reported with the fields that differ; a trade with no candidate is unmatched.
\end{definition}

\omcode{code/firm/posttrade/firm_posttrade.py}{43}{65}{Confirmation matching: candidates by key, the best candidate by fewest differing fields,
and a partial match reported with its score and the fields that differ.}\label{lst:pl:post-trade-systems:match}

The key fields must be exact, because they decide which record is compared with which; the tolerances apply to what is compared. Quantities,
dates and accounts have no tolerance: a trade for 9\,500 shares is not a trade for 9\,400. The price is the one field where a tolerance makes
sense, because the firm and the broker compute an average price over many fills and may carry it to different precision.

\begin{omfigure}
\begin{tikzpicture}
\begin{semilogxaxis}[width=11.5cm,height=5.6cm,xlabel={price tolerance (\$ per share, log scale)},ylabel={matched automatically (\%)},
  grid=major,grid style={gray!20},tick label style={font=\small},label style={font=\small},table/col sep=comma,
  ymin=55,ymax=104,ytick={60,70,80,90,100},xmin=5e-8,xmax=0.1]
\addplot[omDef,thick,mark=*] table[x=tolerance,y=rate_pct]{figdata/platforms/22-post-trade-systems/match_rate.csv};
\node[omDef,font=\scriptsize,anchor=south] at (axis cs:1e-7,60.4) {60.4};
\node[omDef,font=\scriptsize,anchor=south] at (axis cs:1e-6,61.0) {61.0};
\node[omDef,font=\scriptsize,anchor=south east] at (axis cs:1e-5,67.6) {67.6};
\node[omDef,font=\scriptsize,anchor=south] at (axis cs:1e-3,95.8) {95.8};
\node[omDef,font=\scriptsize,anchor=south] at (axis cs:5e-2,96.8) {96.8};
\draw[omThm,dashed] (axis cs:5e-5,55) -- (axis cs:5e-5,100) node[pos=0.15,right,font=\scriptsize,align=left]{half a unit in\\the fourth decimal};
\end{semilogxaxis}
\end{tikzpicture}

\omcaption{Share of Monday's 500 trades matched automatically against the price tolerance. Below \$0.00005 the three brokers that round the
average price to four decimals break on almost every trade (60.4\% matched at \$0.0000001); from \$0.00005 to \$0.01 only the planted
discrepancies remain (95.8\%); at \$0.05 five wrong prices are accepted as matches (96.8\%). Data:
\texttt{fig\_posttrade.py}.}\label{fig:pl:post-trade-systems:tolerance}
\end{omfigure}

\Cref{fig:pl:post-trade-systems:tolerance} shows the two ways to get a tolerance wrong. Too tight, and the rounding convention of three brokers
becomes 177 breaks on Monday that nobody needs to fix: a strict match rate of 60.4\%. Too loose, and a wrong price of a few cents is matched, and
the error is found --- if at all --- when the cash does not agree. Between the two lies a plateau from half a unit in the fourth decimal to a
cent, where the match rate is 95.8\% and every remaining break is a real one. The right tolerance is the plateau's lower end, derived from the
precision each counterparty is known to use, and recorded per counterparty in reference data rather than chosen globally.

\section{Settlement instructions}

Once affirmed, each trade becomes an instruction: for Monday's first trade, a sale of 9\,400 shares at \$97.586941, a \emph{deliver versus
payment} of \$917\,317.25 on the settlement date, from the fund's custody account to the broker's agent and account taken from the SSI table.
Nothing in the instruction is typed; every field comes from the canonical trade or from reference data. That is the point of the SSI table: a
wrong account on one confirmation is a break on one trade, but a wrong account in the SSI table is a fail on every trade with that
counterparty in that market until someone notices, so SSI changes go through the same approval as the parameter changes of
chapter~\ref{ch:pl:signal-serving-and-parameter-management} --- proposed, checked by a second person, effective from a stated date, never
edited in place.

In the model, wrong accounts are the most common discrepancy: 43 of the week's 138 breaks, against 32 missing confirmations, 27 wrong
settlement dates, 26 wrong prices and 10 wrong quantities (\cref{fig:pl:post-trade-systems:kinds}). They are also the kind a firm can remove at
the source, by maintaining SSIs once and sharing them with its counterparties instead of letting each side key them.

\section{Reconciliation and breaks}

\begin{definition}[Reconciliation, reconciliation break]\label{def:pl:post-trade-systems:recon}
\emph{Reconciliation}\index{reconciliation} is the periodic comparison of two records that should agree --- the firm's positions or cash
against a custodian's or prime broker's statement, the firm's trades against a clearing broker's --- item by item. A \emph{reconciliation
break}\index{reconciliation break} is an item on which they differ: present on one side only, or present on both with different values.
\end{definition}

\omcode{code/firm/posttrade/firm_posttrade.py}{106}{120}{Reconciliation by key, with one-to-many groups: a custodian's block position is
compared with the sum of the firm's allocations to it.}\label{lst:pl:post-trade-systems:recon}

The difficulty of \omterm{def:pl:post-trade-systems:recon}{reconciliation} is rarely the comparison; it is the keys. The custodian may report one position where the firm holds five
allocations (a one-to-many group), a cash movement may be the sum of a day's trades (an aggregate), and identifiers may differ between the two
sides, so the reference data of chapter~\ref{ch:pl:reference-data-and-symbology-services} decides whether a difference is real. A break then
needs three more things before it is useful: a kind (missing on one side, a quantity difference), an owner, and an age.

\begin{definition}[Break ageing]\label{def:pl:post-trade-systems:ageing}
\emph{Break ageing}\index{break ageing} is the classification of open breaks by the time since they were opened, in buckets (under a day, one
to three days, three to seven, more), reported daily with the owner of each break; the old buckets are where risk hides, because a break that
survives several \omterm{def:pl:post-trade-systems:recon}{reconciliations} is one that nobody understands.
\end{definition}

\section{Operating under a short settlement cycle}

The model's operations team is three people. Breaks are worked largest value first, each taking a random time whose average depends on the kind
(30~minutes for a settlement date, 45 for a price or a quantity, an hour for an account, an hour and a half to chase a missing confirmation). The
window is where the two cycles differ. Under T+2 the model gives the team the next working day, 08:00 to 17:00, before a trade can no longer be
settled on time; under T+1 it gives the evening of the trade date, 17:00 to the 21:00 \omterm{def:pl:post-trade-systems:confirmation}{affirmation} cut-off. These windows are the model's
choices, set on the cut-offs quoted in the hook; a real team also works during the trading day, which helps both cycles equally.

\omcode{code/platforms/22-post-trade-systems/python/pl_posttrade.py}{93}{109}{The team works the breaks largest value first; a break that
cannot be finished inside the window fails to settle.}\label{lst:pl:post-trade-systems:resolve}

\begin{table}[htbp]
\centering\small
\begin{tabular}{lrrrrrrr}
\toprule
cycle & breaks & fixed & failed & value failed & recon breaks & aged 1--3 days & aged 3--7 days\\
\midrule
T+2 & 138 & 128 & 10 & \$2.74m & 4 & 4 & 0\\
T+1 & 138 & 91 & 47 & \$19.17m & 24 & 12 & 5\\
\bottomrule
\end{tabular}
\caption{The same week of 2\,500 trades and 138 breaks under the two cycles, with three people in operations. \omterm{def:pl:post-trade-systems:recon}{Reconciliation} with the custodian
on Friday evening covers the trades due to settle by Friday; under T+1, 7 of the 24 breaks are under a day old.}
\label{tab:pl:post-trade-systems:cycles}
\end{table}

The breaks are the same under both cycles --- they come from the confirmations, not the calendar --- but the work they need does not fit in the
evening. The week's breaks need about 26 hours of work a day on average; the next working day gives three people 27 hours, the evening 12. Under
T+2 ten trades fail in the week, \$2.74~million of value; under T+1 forty-seven fail, \$19.17~million, and the Friday \omterm{def:pl:post-trade-systems:recon}{reconciliation} finds six
times as many position breaks (\cref{tab:pl:post-trade-systems:cycles}). A failed trade is not lost --- it settles later --- but the fund does
not have the securities or the cash it expected, the \omterm{def:pl:position-and-pnl-services:service}{position service} of chapter~\ref{ch:pl:position-and-pnl-services} must carry settled and
unsettled positions separately, and some markets charge for the fail: Book~2, chapter~5, describes the fails charge in the Treasury market.

\begin{omfigure}
\begin{tikzpicture}
\begin{axis}[width=11.5cm,height=5.4cm,ybar,bar width=9pt,ylabel={trades},symbolic x coords={account,missing,settle_date,price,qty},
  xticklabels={wrong\\account,missing,wrong\\settlement date,wrong\\price,wrong\\quantity},xtick=data,
  xticklabel style={align=center},
  grid=major,grid style={gray!20},tick label style={font=\small},label style={font=\small},table/col sep=comma,ymin=0,ymax=52,
  nodes near coords,every node near coord/.append style={font=\scriptsize},
  legend style={font=\small,at={(0.5,1.03)},anchor=south,legend columns=3,draw=none,
    /tikz/every even column/.append style={column sep=4mm}}]
\addplot[fill=omDef!35,draw=omDef] table[x=kind,y=breaks]{figdata/platforms/22-post-trade-systems/kinds.csv};
\addplot[fill=omThm!25,draw=omThm] table[x=kind,y=failed_t2]{figdata/platforms/22-post-trade-systems/kinds.csv};
\addplot[fill=omThm!70,draw=omThm] table[x=kind,y=failed_t1]{figdata/platforms/22-post-trade-systems/kinds.csv};
\legend{breaks in the week,failed under T+2,failed under T+1}
\end{axis}
\end{tikzpicture}

\omcaption{The week's 138 breaks by kind, and the trades that failed to settle under each cycle. Wrong accounts and missing confirmations are
the most frequent and the slowest to fix, and together make 36 of the 47 fails under T+1. Data: \texttt{fig\_posttrade.py}.}
\label{fig:pl:post-trade-systems:kinds}
\end{omfigure}

There are two ways back to the fails of the longer cycle, and \cref{fig:pl:post-trade-systems:staffing} compares them. The first is people:
with seven in operations instead of three the evening absorbs the work, and the week's fails fall to 11. The second is fewer breaks: halving
the rate of wrong accounts and missing confirmations --- SSIs maintained once and shared, allocations sent to the broker as soon as they are
decided --- brings the fails to 13 with the same three people. The second is cheaper and it lasts, which is why the move to T+1 was prepared
as much in reference data and allocation workflows as in staffing.

\begin{omfigure}
\begin{tikzpicture}
\begin{axis}[width=11.5cm,height=5.6cm,xlabel={people in operations},ylabel={trades failed in the week (T+1)},
  grid=major,grid style={gray!20},tick label style={font=\small},label style={font=\small},table/col sep=comma,
  xmin=2.6,xmax=8.6,ymin=0,ymax=52,xtick={3,4,5,6,7,8}]
\addplot[omThm,thick,mark=*] table[x=staff,y=failed]{figdata/platforms/22-post-trade-systems/staffing.csv};
\addplot[omDef,thick,mark=square*] table[x=staff,y=failed_half_errors]{figdata/platforms/22-post-trade-systems/staffing.csv};
\draw[gray,dashed] (axis cs:2.6,10) -- (axis cs:8.6,10) node[pos=0.97,below left,font=\scriptsize]{T+2, three people: 10};
\node[font=\scriptsize,anchor=north east,align=right] at (axis cs:8.5,50) {\textcolor{omThm}{$\bullet$ the model's error rates}\\
  \textcolor{omDef}{$\blacksquare$ wrong accounts and missing confirmations halved}};
\end{axis}
\end{tikzpicture}

\omcaption{Trades failed in the week under T+1 against the size of the operations team, at the model's error rates (red) and with the
wrong accounts and missing confirmations halved (blue). Seven people bring the fails to 11; halving the two commonest errors brings them to 13
with three. Data: \texttt{fig\_posttrade.py}.}\label{fig:pl:post-trade-systems:staffing}
\end{omfigure}

The general lesson is the one of chapter~\ref{ch:pl:trade-capture-and-booking}: post-trade cost is set upstream. Every break the pipeline
creates is paid for by a person against a deadline, and a shorter cycle shortens the deadline without changing the number of breaks. The
measure to watch is the share of trades that go from booking to settlement untouched --- the straight-through rate --- by counterparty and by
kind of break, together with the age of what is left.

\begin{dated}{2026-09}{The US move to T+1}\label{dat:pl:post-trade-systems:t1}
SEC Release 34-96930 (2023) amended Rule 15c6-1 to shorten the standard settlement cycle for most broker-dealer transactions from T+2 to T+1,
with a compliance date of 28 May 2024, the day after a Federal holiday. New Rule 15c6-2 requires brokers to have agreements or policies aimed
at completing allocations, confirmations and \omterm{def:pl:post-trade-systems:confirmation}{affirmations} as soon as technologically practicable and no later than the end of the trade date;
new Rule 17Ad-27 requires clearing agencies that provide a central matching service to facilitate \omterm{def:pl:trade-capture-and-booking:stp}{straight-through processing} and report on it
each year. The release cites the industry's T+1 Report, which moved the \omterm{def:pl:post-trade-systems:confirmation}{affirmation} cut-off from 11:30~a.m.\ ET on the day after the trade to
9:00~p.m.\ ET on the trade date.
\end{dated}

\section{Tutorial: a week of confirmations under two cycles}\label{tut:pl:post-trade-systems:tut}

\begin{tutorial}
\textbf{Goal.} Match a week of trades with planted discrepancies, instruct, reconcile and age the breaks, then replay the week under T+1.
\textbf{End state:} \cref{fig:pl:post-trade-systems:tolerance} and \cref{tab:pl:post-trade-systems:cycles}.
\begin{tutsteps}
\item \textbf{Trades and confirmations}: \texttt{pl\_posttrade.trades(0, 1)} and \texttt{confirms(ours)}; count the planted kinds.
\item \textbf{Matching}: \texttt{match\_rates(tols=\dots)} over tolerances from $10^{-7}$ to $0.05$.
\item \textbf{Instructions}: build an \texttt{SSITable} for the eight brokers and generate the instruction of every matched trade.
\item \textbf{Breaks and fixes}: \texttt{week(\textquotedbl T+2\textquotedbl)} and \texttt{week(\textquotedbl T+1\textquotedbl)}.
\item \textbf{\omterm{def:pl:post-trade-systems:recon}{Reconciliation}}: \texttt{custodian\_recon(cycle)}, the breaks and their ageing.
\end{tutsteps}
\textbf{What to change next.} Add a cash \omterm{def:pl:post-trade-systems:recon}{reconciliation} (the day's net payments per fund against the custodian's cash statement) with
aggregate matching; let the team work part of the trading day and see whether T+1 still needs more people.
\end{tutorial}

\section{Build: post-trade}\label{bld:pl:post-trade-systems:posttrade}

\begin{build}
\bfield{Purpose} Agree every trade with its counterparty, instruct it, and prove afterwards that the books agree with the custodian, with every
exception in one queue.
\bfield{Interface} \texttt{match}, \texttt{match\_rate}, \texttt{MatchResult}, \texttt{SSITable.instruction}, \texttt{reconcile} (with
\texttt{groups}), \texttt{Break}, \texttt{classify}, \texttt{age}.
\bfield{Rules} Keys exact, tolerances per field and recorded; instructions generated from SSIs, never typed; SSI changes approved and dated;
every break has a kind, an owner and an opening time; ageing reported daily.
\bfield{Acceptance tests} \texttt{code/firm/posttrade/tests/}: matched, partial and unmatched trades on both sides; the best candidate chosen
among several; an instruction from the SSI table and a missing SSI; one-to-many \omterm{def:pl:post-trade-systems:recon}{reconciliation} and the ageing buckets.
\bfield{Stretch} Cash and trade \omterm{def:pl:post-trade-systems:recon}{reconciliation} with aggregate matching; an \omterm{def:pl:post-trade-systems:confirmation}{affirmation} workflow with cut-off times and escalation; a
straight-through rate report by counterparty.
\end{build}

\begin{omsources}
\item US Securities and Exchange Commission, \emph{Shortening the Securities Transaction Settlement Cycle}, Release No.~34-96930, 2023.
\item One Quant Book~1, chapter~5 (settlement, the settlement cycle, fails, delivery versus payment); Book~2, chapter~5 (the fails charge);
Book~16, chapter~4 (the custodian).
\end{omsources}

\section{Exercises}

\begin{exercise}[$\star$]\label{exo:pl:post-trade-systems:1}
Why are the key fields matched exactly and only the other fields with a tolerance?
\end{exercise}

\begin{exercise}[$\star$]\label{exo:pl:post-trade-systems:2}
A broker rounds prices to four decimals. What is the smallest price tolerance at which its correct confirmations always match?
\end{exercise}

\begin{exercise}[$\star$]\label{exo:pl:post-trade-systems:3}
Under T+1, how many of the week's breaks were fixed, and what share of the breaks failed?
\end{exercise}

\begin{exercise}[$\star\star$]\label{exo:pl:post-trade-systems:4}
Why is a wrong account in the SSI table worse than a wrong account on one confirmation?
\end{exercise}

\begin{exercise}[$\star\star$]\label{exo:pl:post-trade-systems:5}
The custodian reports one position of 3\,000 shares where the firm holds allocations of 1\,000 and 2\,000 to two funds under the same account.
How does the \omterm{def:pl:post-trade-systems:recon}{reconciliation} avoid three breaks?
\end{exercise}

\begin{exercise}[$\star\star$]\label{exo:pl:post-trade-systems:6}
Why does the Friday \omterm{def:pl:post-trade-systems:recon}{reconciliation} find only 4 position breaks under T+2 when 10 trades failed?
\end{exercise}

\begin{exercise}[$\star\star\star$]\label{exo:pl:post-trade-systems:7}
\emph{Coding.} Rerun \texttt{week(\textquotedbl T+1\textquotedbl, staff=s)} for $s$ from 3 to 8 and find the smallest team that fails no more
trades than T+2 with three people.
\end{exercise}

\begin{exercise}[$\star\star\star$]\label{exo:pl:post-trade-systems:8}
\emph{Find the flaw.} \textquotedbl{}Our match rate went from 95.8\% to 96.8\% this month after we widened the price tolerance to five cents:
operations is doing better.\textquotedbl{}
\end{exercise}

\section{Problem: One Evening Instead of a Day}

\begin{problem}[{Weekend problem --- one evening instead of two days}]\label{pb:pl:post-trade-systems:1}
The chapter's week: 500 trades a day, eight brokers, five funds, planted discrepancies, three people in operations.

\medskip\noindent\textbf{Part I --- Agreeing the trade.}
\begin{enumerate}
\item What is compared between a booked trade and a confirmation, and what must be exact?
\item What does \omterm{def:pl:post-trade-systems:confirmation}{affirmation} add to matching?
\item Why do three brokers produce breaks at a tight tolerance, and how many trades match at \$0.0000001?
\item Where is the plateau of the match rate, and what is the rate there?
\item What happens at a tolerance of five cents?
\end{enumerate}

\medskip\noindent\textbf{Part II --- Instructing and reconciling.}
\begin{enumerate}[resume]
\item What does the \omterm{def:pl:post-trade-systems:ssi}{settlement instruction} of Monday's first trade contain, and where does each field come from?
\item How should SSI changes be controlled?
\item What makes \omterm{def:pl:post-trade-systems:recon}{reconciliation} hard in practice?
\item What does a break need beyond the difference itself?
\item Why are the old ageing buckets the ones to watch?
\end{enumerate}

\medskip\noindent\textbf{Part III --- The two cycles.}
\begin{enumerate}[resume]
\item What window does the model give the team under each cycle, and why?
\item How much work do the breaks need per day, and how much does each window give?
\item How many trades fail under each cycle, and for how much value?
\item Which kinds of break fail most under T+1?
\item What does the Friday \omterm{def:pl:post-trade-systems:recon}{reconciliation} find under each cycle?
\end{enumerate}

\medskip\noindent\textbf{Part IV --- The verdict.}
\begin{enumerate}[resume]
\item State the \emph{named result}: the share of trades matched automatically at each tolerance, the breaks remaining at the \omterm{def:pl:post-trade-systems:confirmation}{affirmation}
cut-off under a two-day and a one-day cycle, and the fails that follow.
\item How many people does T+1 need to fail no more trades than T+2 with three?
\item What does halving wrong accounts and missing confirmations achieve?
\item Which of the two would you fund, and why?
\item In one sentence: what does a shorter settlement cycle change?
\end{enumerate}
\end{problem}

\section{Interview questions}

\begin{interviewq}[$\star$ \iqroles{developer}]\label{iq:pl:post-trade-systems:1}
What is a \omterm{def:pl:post-trade-systems:recon}{reconciliation break}, and what would you put in a break report?
\end{interviewq}

\begin{interviewq}[$\star\star$ \iqroles{developer}]\label{iq:pl:post-trade-systems:2}
How do you choose the tolerances of a confirmation-matching engine?
\end{interviewq}

\begin{interviewq}[$\star\star$ \iqroles{developer}]\label{iq:pl:post-trade-systems:3}
Your custodian reports positions by account and your books hold them by fund and strategy. Design the \omterm{def:pl:post-trade-systems:recon}{reconciliation}.
\end{interviewq}

\begin{interviewq}[$\star\star$ \iqroles{developer}]\label{iq:pl:post-trade-systems:4}
The market moves to a one-day settlement cycle. What changes in your post-trade systems and team?
\end{interviewq}

\begin{interviewq}[$\star\star\star$ \iqroles{developer}]\label{iq:pl:post-trade-systems:5}
Design the post-trade platform of a firm that trades equities with thirty brokers for a hundred funds.
\end{interviewq}
